\documentclass[12pt,a4paper]{article}
\usepackage[T2A]{fontenc}
\usepackage[utf8]{inputenc}
\usepackage[russian]{babel}
\usepackage[a4paper,left=20mm,right=20mm,top=22mm,bottom=22mm]{geometry}
\usepackage{amsmath,amssymb,mathtools}
\usepackage{array,booktabs,longtable}
\usepackage{xcolor}
\usepackage{tikz}
\usepackage[unicode,colorlinks=true,linkcolor=blue!50!black]{hyperref}
\setlength{\parindent}{0pt}
\setlength{\parskip}{7pt}
\setlength{\emergencystretch}{3em}
\newcommand{\gcdru}{\mathop{\text{НОД}}\nolimits}
\newcommand{\lcmru}{\mathop{\text{НОК}}\nolimits}
\newcommand{\subheading}[1]{\par\smallskip\textbf{#1}\par}

\begin{document}
\begin{titlepage}
\centering
\vspace*{0.33\textheight}
{\LARGE\bfseries Ревью домашней работы\par}
\vspace{18mm}
{\large Автор: Лёвин Артём Александрович\par}
\vspace{7mm}
{\large 10 октября 2026 года\par}
\vfill
\begin{tikzpicture}
\draw[blue!40!black,line width=0.65pt]
  (0,0) .. controls (2.0,0.55) and (3.4,-0.55) .. (5.2,0)
  .. controls (6.8,0.42) and (8.4,-0.40) .. (10,0);
\end{tikzpicture}
\vspace*{9mm}
\end{titlepage}

\newpage
\section*{Сводная таблица заданий с замечаниями}
\small
\renewcommand{\arraystretch}{1.5}
\begin{tabular}{@{}p{2.2cm}p{2.5cm}p{5.2cm}p{1.9cm}p{3.2cm}@{}}
\toprule
\textbf{№ задания} & \textbf{Тема} & \textbf{Суть замечания} & \textbf{Ваш ответ} & \textbf{Правильный ответ} \\
\midrule
\hyperref[task:oct09-3]{9 октября, № 3}
& НОД и НОК двух чисел
& Не проверены отдельно оба условия задачи. Проверено только равенство произведений.
& $n=90$
& $n=90$; обе проверки выполнены. \\
\addlinespace
\hyperref[task:oct10-3]{10 октября, № 3}
& Все пары чисел по НОД и НОК
& Найдено произведение чисел, но не перечислены пары и не доказана полнота списка.
& Не указан
& $(18,540)$, $(36,270)$, $(54,180)$, $(90,108)$. \\
\bottomrule
\end{tabular}
\normalsize

\newpage
\thispagestyle{empty}
\null\vfill
\begin{center}
{\Huge\bfseries Разбор ошибок}
\end{center}
\vfill\null

\newpage
\phantomsection\label{task:oct09-3}
\section*{9 октября 2026 года. Задание № 3}

\subheading{Текст задания}
Для натурального $n$ известно, что
\[
\gcdru(n,72)=18,\qquad \lcmru(n,72)=360.
\]
Найдите $n$ и проверьте оба условия.

\subheading{Ваше решение}
Вы правильно использовали свойство произведения НОД и НОК двух натуральных чисел:
\[
\gcdru(n,72)\cdot\lcmru(n,72)=n\cdot72.
\]
Подставили данные задачи и нашли число:
\[
18\cdot360=72n,\qquad
n=\frac{18\cdot360}{72}=90.
\]
После этого Вы записали проверку равенства произведений:
\[
360\cdot18=72\cdot90,\qquad 6480=6480.
\]

\subheading{Ваш ответ}
$n=90$.

\subheading{Ошибка}
\textcolor{red}{Ошибка:} не выполнено обязательное требование проверить \emph{оба} условия задачи. Найденное число верное, однако равенство произведений само по себе не заменяет вычисление НОД и НОК.

\subheading{Где именно возникает ошибка}
Последний правильный шаг --- получение числа $n=90$. Следующий шаг с равенством $360\cdot18=72\cdot90$ тоже арифметически верен. Недочёт возникает \emph{после него}: решение заканчивается, хотя нужно отдельно установить
\[
\gcdru(90,72)=18\quad\text{и}\quad\lcmru(90,72)=360.
\]
Эти два утверждения не были показаны.

\subheading{Почему этого недостаточно}
Формула
\[
\gcdru(a,b)\cdot\lcmru(a,b)=ab
\]
связывает два искомых значения только через их произведение. Из совпадения произведений нельзя заключить, что каждый множитель вычислен правильно: например, произведения $6\cdot36$ и $3\cdot72$ равны, но сами пары множителей различны. Поэтому проверка $6480=6480$ подтверждает арифметику поиска числа, но не подтверждает отдельно указанные НОД и НОК.

\subheading{Ключевая идея}
\textcolor{blue!60!black}{Ключевая идея:} сначала применить связь НОД и НОК для поиска кандидата, затем \emph{непосредственно} вычислить НОД найденного числа и числа $72$, после чего вычислить их НОК.

\subheading{Исправление от последнего правильного шага}
Сохраняем полученное значение $n=90$. Проверяем первое условие алгоритмом Евклида:
\[
90=72+18,\qquad 72=18\cdot4+0.
\]
Последний ненулевой остаток равен $18$, следовательно,
\[
\gcdru(90,72)=18.
\]
Теперь проверяем второе условие, используя уже \emph{вычисленный} НОД:
\[
\lcmru(90,72)=\frac{90\cdot72}{\gcdru(90,72)}
=\frac{90\cdot72}{18}=90\cdot4=360.
\]
Оба условия выполнены.

\subheading{Полное правильное решение}
По условию НОД числа $n$ и числа $72$ равен $18$, а НОК равен $360$. Для любых двух натуральных чисел произведение НОД и НОК равно произведению самих чисел. Поэтому
\[
18\cdot360=n\cdot72.
\]
Разделим обе части равенства на $72$ (это число не равно нулю):
\[
n=\frac{18\cdot360}{72}
=18\cdot5=90.
\]
Число $90$ натуральное, значит, требование к $n$ соблюдено. Но значение $n$ пока только кандидат, поэтому нужно подтвердить исходные условия. По алгоритму Евклида $\gcdru(90,72)=18$. Далее
\[
\lcmru(90,72)=\frac{90\cdot72}{18}=360.
\]
Таким образом, найденное число удовлетворяет двум условиям одновременно. Кроме того, исходное равенство $72n=18\cdot360$ однозначно задаёт $n$, так что другого значения быть не может.

\subheading{Проверка}
\[
\underbrace{\gcdru(90,72)=18}_{\text{первое условие}},
\qquad
\underbrace{\lcmru(90,72)=360}_{\text{второе условие}}.
\]

\subheading{Итог}
\textcolor{teal!60!black}{Итог:} Ваш ответ $n=90$ правильный. Чтобы решение было полным, необходимо отдельно подтвердить оба данных в условии значения: НОД равен $18$, НОК равен $360$.

\subheading{Задача на закрепление}
Для натурального $n$ известно, что $\gcdru(n,60)=12$ и $\lcmru(n,60)=420$. Найдите $n$ и проверьте оба условия.

\newpage
\phantomsection\label{task:oct10-3}
\section*{10 октября 2026 года. Задание № 3}

\subheading{Текст задания}
Найдите \emph{все неупорядоченные пары} натуральных чисел $(a,b)$, для которых
\[
\gcdru(a,b)=18,\qquad\lcmru(a,b)=540.
\]
Обоснуйте, что других пар нет.

\subheading{Ваше решение}
Вы верно записали формулу связи НОД и НОК:
\[
\gcdru(a,b)\cdot\lcmru(a,b)=ab.
\]
Затем получили
\[
18\cdot540=9720,\qquad ab=9720.
\]
Также верно начали разложение данных чисел на простые множители:
\[
18=2\cdot3^2,\qquad 540=2^2\cdot3^3\cdot5.
\]
Однако на этом решение заканчивается: пары $(a,b)$ не перечислены и доказательство того, что других пар нет, отсутствует.

\subheading{Ваш ответ}
Не указан.

\subheading{Ошибка}
\textcolor{red}{Ошибка:} задача не доведена до ответа. Из условия $ab=9720$ ещё не следует, какие именно числа $a$ и $b$ имеют НОД $18$ и НОК $540$. Нужно учесть оба ограничения и перечислить \emph{все} допустимые пары.

\subheading{Где именно возникает ошибка}
Последний правильный шаг --- запись $ab=9720$ и разложений чисел $18$ и $540$ на простые множители. Первый недостающий шаг --- использование условия $\gcdru(a,b)=18$ для представления $a=18x$, $b=18y$, где $x$ и $y$ взаимно просты. Без этого произведение $ab$ допускает много неподходящих пар множителей. Например, числа $18$ и $540$ подходят, а числа $27$ и $360$, хотя их произведение также равно $9720$, не подходят: их НОД равен $9$, а не $18$.

\subheading{Почему это важно}
Условие о произведении является необходимым, но недостаточным. По требованию задачи нужно выполнить сразу \emph{два} условия: получить заданный НОД и заданный НОК. Для этого удобно отделить общий делитель $18$ от каждого числа. После выделения НОД у оставшихся множителей не должно быть общих простых делителей, иначе общий делитель исходных чисел оказался бы больше $18$.

\subheading{Ключевая идея}
\textcolor{blue!60!black}{Ключевая идея:} записать $a=18x$ и $b=18y$, где $\gcdru(x,y)=1$. Тогда все способы получить пару $(a,b)$ сводятся к разбиению простых множителей числа $30$ между взаимно простыми числами $x$ и $y$.

\subheading{Исправление от последнего правильного шага}
Начнём с Вашего верного равенства $ab=9720$ и условия $\gcdru(a,b)=18$. Поскольку $18$ делит оба числа, существуют натуральные $x,y$ такие, что
\[
a=18x,\qquad b=18y,\qquad\gcdru(x,y)=1.
\]
Подставим в найденное произведение:
\[
18x\cdot18y=9720,
\qquad xy=\frac{9720}{18^2}=30.
\]
Число $30$ раскладывается так:
\[
30=2\cdot3\cdot5.
\]
Каждый из трёх различных простых множителей должен целиком принадлежать либо $x$, либо $y$: делить один простой множитель между ними нельзя, иначе он стал бы общим делителем. С учётом того, что порядок чисел в паре не важен, получаем
\[
(x,y)=(1,30),\ (2,15),\ (3,10),\ (5,6).
\]
Умножаем оба числа в каждой паре на $18$:
\[
(a,b)=(18,540),\ (36,270),\ (54,180),\ (90,108).
\]

\subheading{Полное правильное решение}
Так как $\gcdru(a,b)=18$, числа $a$ и $b$ кратны $18$. Обозначим их частные через $x$ и $y$:
\[
a=18x,\quad b=18y.
\]
Условие $\gcdru(a,b)=18$ равносильно тому, что $\gcdru(x,y)=1$. Если бы $x$ и $y$ имели общий делитель больше единицы, то НОД чисел $a$ и $b$ был бы больше $18$.

Для взаимно простых $x,y$ имеем $\lcmru(x,y)=xy$. При умножении обоих чисел на $18$ их НОК также умножается на $18$, следовательно,
\[
\lcmru(a,b)=18\,\lcmru(x,y)=18xy.
\]
По второму условию $18xy=540$, откуда
\[
xy=30.
\]
Все неупорядоченные пары делителей с произведением $30$ имеют вид
\[
(1,30),\quad(2,15),\quad(3,10),\quad(5,6).
\]
Во всех четырёх парах числа взаимно просты. Других неупорядоченных пар натуральных чисел с произведением $30$ нет: меньший множитель не превосходит $\sqrt{30}<6$, поэтому достаточно проверить делители числа $30$, не превышающие $5$, а именно $1,2,3,5$.

Возвращаясь к исходным числам, умножим каждую компоненту пары на $18$. Получаем ровно четыре пары:
\[
\boxed{(18,540),\quad(36,270),\quad(54,180),\quad(90,108).}
\]
Перестановка двух чисел не даёт новую пару, потому что по условию пары неупорядоченные.

\subheading{Проверка}
Для каждой найденной пары разделим оба числа на $18$. Получаем взаимно простые пары $(1,30)$, $(2,15)$, $(3,10)$, $(5,6)$, причём в каждой произведение равно $30$. Значит, у каждой исходной пары НОД равен $18$, а НОК равен $18\cdot30=540$. Проверка обоих условий выполнена для всех четырёх пар.

\subheading{Итог}
\textcolor{teal!60!black}{Итог:} формула $ab=9720$ найдена правильно, но её недостаточно для полного ответа. Всего существует четыре неупорядоченные пары: $(18,540)$, $(36,270)$, $(54,180)$ и $(90,108)$. Других пар нет, поскольку перечислены все допустимые взаимно простые пары множителей числа $30$.

\subheading{Задача на закрепление}
Найдите все неупорядоченные пары натуральных чисел $(a,b)$, для которых $\gcdru(a,b)=12$ и $\lcmru(a,b)=1260$. Обоснуйте, что других пар нет.

\vfill
\begin{center}
\footnotesize Лёвин Артём Александрович эксклюзивно для Ксении Клыковой
\end{center}

\end{document}